% ============================================================
%  Глава 4 — Динамика: инерция моторов
% ============================================================
\chapter{Динамика: моторы не отрабатывают мгновенно}
\label{ch:dynamics}

\section{Вопрос}

В предыдущей главе мы кормили в кинематику $v$ и $\omega$ так, будто
\emph{что задали --- то и едет}. На самом деле это не так. Если в
момент $t=0$ сказать роботу «езжай $0.2$~м/с» (раньше стоял), он
\emph{не сразу} достигнет $0.2$~м/с --- сначала будет разгоняться.

\begin{ideabox}
Между \emph{командой} $u_v$ и \emph{фактической} скоростью $v$ есть
зазор: моторам нужно время, чтобы раскрутить инерцию, преодолеть
противо-ЭДС и трение. Длительность этого зазора --- постоянная времени
$\tau_v$.

Это \textbf{динамика} робота. Без неё модель неполна, и MPC,
обученный на одной только кинематике, будет систематически ошибаться.
\end{ideabox}

\section{Апериодическое звено первого порядка}

Простейшая модель, которая адекватно описывает разгон моторов под
такой нагрузкой --- \textbf{апериодическое звено первого порядка}
(в литературе также «инерционное звено», PT$_1$):

\begin{formulabox}[title={Динамика линейной скорости}]
\begin{equation}
  \dot{v} = \frac{1}{\tau_v}\,(u_v - v).
  \label{eq:motor_lin_dyn}
\end{equation}
\end{formulabox}

Если $u_v$ постоянно, решение этого уравнения --- классическая экспонента:

\begin{equation}
  v(t) = u_v - (u_v - v_0)\,e^{-t/\tau_v},
  \label{eq:exp_response}
\end{equation}

где $v_0 = v(0)$. За время $t = \tau_v$ робот покрывает $\approx 63\%$
пути от $v_0$ до $u_v$; за $t = 3\tau_v$ --- $\approx 95\%$.

\begin{ideabox}
Можно прочитать \eqref{eq:motor_lin_dyn} «по-человечески»: \emph{скорость
изменения скорости пропорциональна тому, как сильно она отстаёт от
команды}. Чем больше отставание --- тем сильнее «толкаем»; когда $v$
догнала $u_v$ --- толкать перестаём.

Аналогия: вода, текущая в бак, где скорость наполнения пропорциональна
разности уровней. RC-цепь в электронике --- та же самая динамика, потому
её и называют «RC-звено».
\end{ideabox}

\begin{examplebox}[title={Численный пример}]
Пусть $\tau_v = 0.15$~с (типичное значение для шасси SAMURAI).
Робот стоит ($v_0 = 0$), команда $u_v = 0.2$~м/с.

\begin{tabular}{rrr}
\toprule
$t$, с & $v$, м/с & доля от $u_v$\\
\midrule
$0.05$ & $0.2 \cdot (1 - e^{-1/3}) \approx 0.057$ & 28.3\%\\
$0.15$ & $0.2 \cdot (1 - e^{-1}) \approx 0.126$ & 63.2\%\\
$0.30$ & $0.2 \cdot (1 - e^{-2}) \approx 0.173$ & 86.5\%\\
$0.45$ & $0.2 \cdot (1 - e^{-3}) \approx 0.190$ & 95.0\%\\
\bottomrule
\end{tabular}

За $T_s = 0.05$~с (один тик MPC) робот успевает разогнаться только
до 28\% желаемого. Это и есть то, что MPC обязан знать.
\end{examplebox}

\section{Аналогично --- угловая скорость}

\begin{formulabox}[title={Динамика угловой скорости}]
\begin{equation}
  \dot{\omega} = \frac{1}{\tau_\omega}\,(u_\omega - \omega).
  \label{eq:motor_ang_dyn}
\end{equation}
\end{formulabox}

Постоянная $\tau_\omega \approx 0.10$~с заметно \emph{меньше} $\tau_v$:
крутить шасси на месте проще, чем разгонять центр массы.

\section{Полная нелинейная ОДУ-модель}

Соберём кинематику \eqref{eq:world_kinematics} и динамику
\eqref{eq:motor_lin_dyn}, \eqref{eq:motor_ang_dyn} в один вектор.
Возьмём «физическое» (легаси) состояние из главы~\ref{ch:state}:

\[
  \vect{x}(t) = \big(p_x,\; p_y,\; \theta,\; v,\; \omega\big)\T \in \Real^5,
  \qquad
  \vect{u}(t) = (u_v,\; u_\omega)\T \in \Real^2.
\]

\begin{formulabox}[title={Нелинейная ОДУ-модель робота}]
\begin{equation}
  \boxed{\;
    \dot{\vect{x}}(t) = \vect{f}(\vect{x}, \vect{u})
    \;=\;
    \begin{pmatrix}
      v\cos\theta \\[3pt]
      v\sin\theta \\[3pt]
      \omega \\[3pt]
      \frac{1}{\tau_v}\,(u_v - v) \\[5pt]
      \frac{1}{\tau_\omega}\,(u_\omega - \omega)
    \end{pmatrix}.\;}
  \label{eq:nonlinear_ode}
\end{equation}
\end{formulabox}

Здесь:
\begin{itemize}
  \item Первые три строки --- \emph{кинематика}. Нелинейность через $\cos, \sin$.
  \item Последние две --- \emph{динамика моторов}. \emph{Линейная}: правая
        часть линейна по $v, \omega, u_v, u_\omega$.
\end{itemize}

\textbf{Вся нелинейность сосредоточена в куске «как $v$ и $\theta$
порождают движение в мире»}. Это важно: при линеаризации (следующая
глава) именно этот кусок упростится.

\section{Откуда берутся $\tau_v$ и $\tau_\omega$}

Постоянные времени --- \emph{физические} параметры, оценённые
экспериментально. В \filepath{tests/test\_motor\_lag.py} есть тест, который
гонит робот на серии импульсов $u_v$ и подбирает $\tau_v$
по методу наименьших квадратов из закона \eqref{eq:exp_response}.
В \filepath{config.yaml} они зашиты как:

\begin{codebox}[title={\filepath{config.yaml}, секция \texttt{control.plant}}]
\begin{verbatim}
control:
  plant:
    Ts:     0.05     # период дискретизации, с
    v0:     0.20     # опорная скорость для линеаризации, м/с
    tau_v:  0.15     # постоянная времени линейной скорости, с
    tau_w:  0.10     # постоянная времени угловой скорости, с
\end{verbatim}
\end{codebox}

\begin{whatifbox}
\textbf{Что если у моторов больше инерция?} Тогда $\tau_v$ увеличится,
и в дискретной модели $1/\tau_v$ уменьшится --- собственное значение
$\lambda = -1/\tau_v$ ближе к нулю. В дискретном виде $|\tilde{\lambda}|
= e^{-T_s/\tau_v}$ ближе к единице --- мотор «медленнее забывает» прошлую
команду. MPC должен это учесть, выбирая более «нежную» $\mat{R}$.

Если поставить в config \emph{неправильно большое} $\tau_v$, ничего
не сломается --- модель просто будет «думать, что моторы вялые». MPC
будет давать более резкие команды (чтоб компенсировать выдуманную
инерцию) --- робот будет дёргаться. Хороший пример «учебной катастрофы»
из задачи курсовой.
\end{whatifbox}

\section{Где это в коде}

\begin{codebox}[title={\filepath{pi\_nodes/control/state\_space\_model.py:29}}]
\begin{lstlisting}[language=Python]
def _continuous_AB(v0: float, tau_v: float, tau_w: float):
    """Build linearised A, B around forward motion at v0."""
    A = np.array([
        [0.0, 0.0, 0.0,        1.0,         0.0       ],
        [0.0, 0.0, v0,         0.0,         0.0       ],
        [0.0, 0.0, 0.0,        0.0,         1.0       ],
        [0.0, 0.0, 0.0,       -1.0/tau_v,   0.0       ],
        [0.0, 0.0, 0.0,        0.0,        -1.0/tau_w ],
    ])
    B = np.array([
        [0.0,        0.0       ],
        [0.0,        0.0       ],
        [0.0,        0.0       ],
        [1.0/tau_v,  0.0       ],
        [0.0,        1.0/tau_w ],
    ])
    return A, B
\end{lstlisting}
\end{codebox}

Это уже \emph{линеаризованная} версия модели --- но мы её ещё не вывели.
Чтобы её получить, нужно сделать линеаризацию \eqref{eq:nonlinear_ode} вокруг
опорной точки. Этим займёмся в следующей главе.

\section{Что мы получили}

\begin{itemize}
  \item Моторы --- \emph{инерционные}: команда отрабатывается за $\tau_v
        \sim 0.15$~с, $\tau_\omega \sim 0.10$~с.
  \item Динамика моторов --- линейная: $\dot{v} = (u_v - v)/\tau_v$ и
        аналогично для $\omega$.
  \item Полная модель робота \eqref{eq:nonlinear_ode} --- пятимерная,
        нелинейная (только в кинематической части), с двумя управлениями.
  \item Постоянные времени --- из тестов на реальном железе, сидят
        в \filepath{config.yaml}.
\end{itemize}

Дальше эту нелинейную систему надо превратить в линейную (для MPC и
ЛКР это базовая необходимость). Идёмм.
